\chapter{Utilities}
\label{chapter:utilities}

This chapter documents reusable utility components that support common development needs across the project. All utility modules are defined in CMake using the naming convention \texttt{\textit{Name}\_UTIL}. These modules can be linked into your targets via CMake's \texttt{target\_link\_libraries} command. This allows modular inclusion of specific utility functionalities based on your build requirements. To use a utility, simply add the corresponding utility target to your CMake target, for example:
\begin{lstlisting}[style=cppstyle]
target_link_libraries(MyTarget PRIVATE Types_UTIL Math_UTIL)
\end{lstlisting}
Below is a list of available utility modules.
\begin{itemize}\itemsep0em
	\item \textbf{Types\_UTIL} – This contains utilities related to type manipulation and handling.
	\item \textbf{Math\_UTIL} – This contains mathematical utilities.
	\item \textbf{System\_UTIL} – This contains utilities related to system commands.
	\item \textbf{Files\_UTIL} – This contains utilities related to file handling.
	\item \textbf{Audio\_UTIL} – This contains utilities related to audio handling.
	\item \textbf{Encryption\_UTIL} – This contains various encryption utilities.
	\item \textbf{ML\_UTIL} – This contains various Machine-learning utilities.
	\item \textbf{Python\_UTIL} – This contains utilities for loading python modules and calling python methods from the c++ code.
	\item \textbf{UI\_UTIL} – This contains user interface input utilities: event types, a shared event queue, and key listeners.
\end{itemize}









\section{Markov Model Generation Module}

These features are all part of the \texttt{ML\_UTIL}.

\subsection{Overview}

In a previous project I was working on for Dungeons \& Dragons, I created a model for generating random names and character sequences. That project code can be found here: 
\begin{center}
\url{https://github.com/torodean/DnD/blob/main/templates/creator.py}. 
\end{center}
The functionality was based on transition patterns between characters of preexisting names. I only later found out that this was referred to as a Markov model. The features related to this which are being added to MIA are a more generalized version of that work, which will follow the naming convention relating them to Markov Models.

The \texttt{MarkovModels.hpp} header defines a templated C++ module for constructing and manipulating first-order discrete-time Markov models of arbitrary order. The implementation is generic and supports any comparable data type, including characters, strings, or custom state representations, etc.

\subsection{First-Order Markov Model}

This section contains an explanation of how the markov model functions. Consider some set of sequences $S$ for some arbitrary type, where each element is denoted by a capitalized letter, $S=\{ABC, ABD, BAD\}$. The probability matrix $P$ is constructed by determining all of the elements which follow another, and at what probability that element has of following the others (The probabilities of elements following some element $K$ is $P_K$). that is $P(S) = \{K:P_K \forall K \in S\}$.

The set of elements which exist in $S$ are ${A, B, C, D, \emptyset}$, where $\emptyset$ denotes the absence of an element (or beginning/end of a sequence). Starting with $A$, we can see that the $A$ element is followed only by $B$ (twice), and $D$ (once) in $S$. The total number of elements ever following an $A$ is thus three. The probabilities following an element $A$ is thus
\begin{align}
P_A = \begin{cases}
  B & :\text{twice} \\
  D & :\text{once}
\end{cases} \implies P_A = \begin{cases}
  B & : 66.\overline{6}\% \\
  D & : 33.\overline{3}\%
\end{cases} = \{B:0.\overline{6}, D:0.\overline{3}\}
\end{align}
Following this same process for the other elements gives
\begin{align}
P_B &= \{A:0.\overline{3}, C:0.\overline{3}, D:0.\overline{3}\} \\ P_C &= P_D = \{\emptyset:1.0\}\\ P_\emptyset &= \{A:0.\overline{6}, B:0.\overline{3}\}.
\end{align}
The total probability matrix for this set of sequences would then be
\begin{align}
P(S) = \{A:P_A, B:P_B, C:P_C, D:P_D, \emptyset: P_\emptyset\} = \begin{cases}
A:\{B:0.\overline{6}, D:0.\overline{3}\} \\
B: \{A:0.\overline{3}, C:0.\overline{3}, D:0.\overline{3}\} \\
C: \{\emptyset:1.0\} \\
D: \{\emptyset:1.0\} \\
\emptyset: \{A:0.\overline{6}, B:0.\overline{3}\}.
\end{cases}
\end{align}
The $\emptyset$ is a special case in that it represents the first character of a sequence (there is never a character after the last). This format may not look like a matrix at all, but it can be re-written to matrix format. First, note that there are a total of 5 elements ($A$, $B$, $C$, $D$, $\emptyset$) which will give a $5 \times 5$ matrix for all possible combinations. The matrix is configured such that both the rows and columns span from $A\rightarrow\emptyset$, covering all the elements of the set. The matrix value of $a, b$ then represents the probability that element $a$ will be proceeded by element $b$.
\begin{align}
P(S) = \left[
\begin{matrix}
0 & 0.\overline{3} & 0 & 0 & 0.\overline{6} \\ 
0.\overline{6} & 0 & 0 & 0 & 0.\overline{3} \\ 
0 & 0.\overline{3} & 0 & 0 & 0 \\ 
0.\overline{3} & 0.\overline{3} & 0 & 0 & 0\\ 
0 & 0 & 1.0 & 1.0 & 0 \\ 
\end{matrix}\right]
\end{align}
This probability matrix thus represents the probability of an element proceeding another in one of the given sequences. It can be used to generate new sequences which adhere to similar patterns of the input sequences. With larger data sets, more possibilities of sequences typically arise as probable outputs. 

One important feature of these models is that under low-entropy (the model is derived from a deterministic source), a uniquely resolvable input set (You can reconstruct exactly one input set) and with enough metadata (initial state, model size, model order, etc), the model can be used to reconstruct the original data.


\subsection{Functionality}

This module provides the following key features:

\begin{itemize}\itemsep0em
    \item \textbf{Probability Matrix Construction:} 
    Generates a normalized transition matrix from a collection of input sequences, where transitions between adjacent elements are counted and then normalized into probabilities.

    \item \textbf{State Querying:}
    Provides utility functions to:
    \begin{itemize}\itemsep0em
        \item Retrieve successor probabilities for a given state.
        \item Check whether a given state exists in the matrix.
        \item Extract the list of all states in the model.
    \end{itemize}

    \item \textbf{Matrix Utilities:}
    Includes functions to:
    \begin{itemize}\itemsep0em
        \item Clear the matrix.
        \item Print the matrix to an output stream.
    \end{itemize}

    \item \textbf{Sampling Support (Planned):}
    A placeholder exists for a future function that will enable sampling the next state from a given current state based on its transition distribution.
\end{itemize}

\subsection{Design Notes}

The transition matrix is represented using nested \texttt{std::unordered\_map} containers. Specifically, the outer map keys are current states, and the values are maps from successor states to transition probabilities. The code leverages C++17 features, including structured bindings and template aliasing, and is contained within the \texttt{markov\_models} namespace for modularity and clarity.



















\section{D0sag3 Command (D3C) Integration} \label{D3C}

These features are all part of the \texttt{Encryption\_UTIL}.

\subsection{D3C Introduction and Overview}

The D3C encryption was an incorporation of an old encryption program I created many years ago as part of the D3C (d0sag3 command) program. The original code was made when I was first learning C++ and this was used as a project for educational purposes. The encryption algorithm utilizes random numbers, bit analysis, variable type conversions, and more. 

\subsection{d0s1 Encryption}

The d0s1 encryption algorithm was the first implementation of encryption within the D3C program. The d0s1 algorithm is programmed solely to encrypt an input string value. To outline the algorithm that d0s1 uses, we will start with an example string "hello." The algorithm follows.

\begin{lstlisting}
	# Start with an input string
	Hello
	
	# Each character is examined individually.
	H e l l o
	
	# The string get's converted to integers based on the ascii value of each character.
	72 101 108 108 111
	
	# The integers are converted to a binary representation.
	1001000 1100101 1101100 1101100 1101111
	
	# A random number is generated for each character that existed.
	103 70 105 115 81
	
	# The random numbers are converted to binary representations.
	1100111 1000110 1101001 1110011 1010001
	
	# The string and random binary numbers are added to a trinary number.
	1001000 1100101 1101100 1101100 1101111
	+	1100111 1000110 1101001 1110011 1010001
	-------------------------------------------
	=	2101111 2100211 2202101 2211111 2111112
	
	# The random numbers selected before are converted to base 12 numbers.
	103 70 105 115 81 = 87 5A 89 97 69
	
	# The base 12 random numbers are placed at the end of the trinary string.
	210111187 21002115A 220210189 221111197 211111269
	
	# The ouput of the encryption is then these values.
	21011118721002115A220210189221111197211111269
\end{lstlisting}

The encryption was meant to have a final stage to decrease the length of the output by assigning different characters to the number sequences output; however, this was never finished.

\subsection{d0s2 Encryption}

d0s2 encryption is a very similar algorithm to d0s1 with one major difference. The encryption of d0s2 requires a user input password that is added into the encryption process. The password and string are both encrypted and then added together in a way that the password is needed for quick decryption.

\subsection{d0s3 Encryption}

The d0s3 encryption algorithm was (as of the time writing this) never finished. The d0s3 was the actual D3C encryption that was originally desired with d0s1 and d0s2 being practice runs for the creator to experiment with C++ first before employing an actual complicated algorithm. MIA currently has parts of the d0s3 encryption programmed in but they are still in development and not yet deployed. The d0s3 encryption algorithm is not related to d0s1 and d0s3 but will instead have a unique and complicated algorithm that can encrypt entire files instead of just string values. The D3C encryption utilities, including d0s1, d0s2, and the partial implementation of d0s3, are included in MIA as general-purpose utilities. These components are designed to be accessible for use across various MIA applications and libraries, providing a consistent and centralized encryption interface where needed.


















\section{Terminal Color Codes}

The terminal color features are part of the \texttt{System\_UTIL} utility. It provides a set of inline functions that insert ANSI escape sequences into output streams to enable colored terminal text. To use it in an application or library, the utility needs included in Cmake, similar to the following example:
\begin{lstlisting}[style=shellstyle]
target_link_libraries(MIATemplate PRIVATE Framework_CORE System_UTIL )
\end{lstlisting}

\subsection*{Usage}

Each color function takes a reference to an \texttt{std::ostream} and returns the stream with the appropriate ANSI code inserted. The \texttt{TerminalColors.hpp} file needs included to use these features. For example, to print red text:

\begin{lstlisting}[style=cppstyle]
#include "TerminalColors.hpp"
	
std::cout << red << "This text is red" << reset << std::endl;
\end{lstlisting}

\subsection*{Available Colors}

\begin{multicols}{2}
\begin{itemize}\itemsep0em
	\item \texttt{red}
	\item \texttt{bgred} (background color)
	\item \texttt{green}
	\item \texttt{bggreen} (background color)
	\item \texttt{yellow}
	\item \texttt{blue}
	\item \texttt{bgblue} (background color)
	\item \texttt{magenta}
	\item \texttt{cyan}
	\item \texttt{white}
	\item \texttt{def} (default color)
	\item \texttt{bgdef} (default background color)
	\item \texttt{reset} (resets all formatting)
\end{itemize}
\end{multicols}
These functions leverage ANSI escape codes and are compatible with Unix-like systems and modern Windows terminals supporting ANSI sequences.


















\section{Virtual Key Strokes Utility}
\label{sec:virtualkeystrokes}

These features are all part of the \texttt{System\_UTIL}.


The \texttt{VirtualKeyStrokes} class provides a cross-platform interface for simulating keyboard and mouse input on Windows and Linux systems. It supports character and string input emulation, basic automation, and platform-specific input features. This utility is designed for automation and testing scenarios requiring simulated user input. It abstracts platform-specific details, providing a unified API for key and mouse event generation.

\textit{Note:} The current implementation is very old code, which could definitely use improvements and may require further work for full feature completeness and robustness.

\subsection{Features}
\begin{itemize}\itemsep0em
	\item Simulate key presses for individual characters and strings.
	\item Support for mouse clicks: left, right, and middle click.
	\item Specialized key presses on Windows (numbers, letters, special keys).
	\item Functions to move the mouse cursor and retrieve pixel color at or near the mouse.
	\item Automation support with sequential key typing and custom delays.
	\item Cross-platform support via Windows API and X11/\texttt{xdo} library on Linux.
\end{itemize}

\subsection{Usage Example}
To simulate typing a string and a mouse left click:
\begin{lstlisting}[style=cppstyle]
virtual_keys::VirtualKeyStrokes vk;
vk.type("Hello, World!");
vk.leftclick();
\end{lstlisting}

\subsection{API Highlights}

\begin{description}
	\item[\texttt{void press(const char\& character, int holdTime = 0, bool verboseMode = false)}]  
	Simulates pressing a single key with optional hold time and verbose output.
	
	\item[\texttt{void type(const std::string\& word)}]  
	Simulates typing a sequence of characters.
	
	\item[\texttt{void mouseClick(ClickType clickType, bool verboseMode = false)}]  
	Simulates mouse clicks of different types (\texttt{LEFT\_CLICK}, \texttt{RIGHT\_CLICK}, \texttt{MIDDLE\_CLICK}).
	
	\item[\texttt{void moveMouseTo(int x, int y)}]  
	Moves the mouse pointer to specified screen coordinates.
	
	\item[\texttt{void getPixelColorAtMouse()}]  
	Retrieves and prints the RGB color of the pixel under the mouse cursor (Windows only).
\end{description}

\subsection{Platform Details}

\begin{itemize}\itemsep0em
	\item \textbf{Windows:} Uses the Windows \texttt{SendInput} API for keyboard and mouse simulation.

	\item \textbf{Linux:} Utilizes the \texttt{xdo} library along with X11 functions to emulate input events.

	\item \textbf{Linux (optional):} The \texttt{xdo} and X11 libraries are optional: when they are not found, the key stroke features are not built and \texttt{System\_UTIL} contains only the other system utilities (see section \ref{section:optional_features}).
\end{itemize}










\section{Type Manipulation Utility}
\label{section:type_utility}

These features are all part of the \texttt{Types\_UTIL}.


\subsection{String Utilities}\label{section:string_utils}

The MIA project contains many helper methods particularly related to string manipulation, parsing, transformation, and formatting. The API for this utility is in the \texttt{StringUtils.hpp} file, where the various helper methods are stored. This file contains the complete documentation for each available function. To use the string utilities, simply include the \texttt{StringUtils.hpp} file and link to the \texttt{Types\_UTIL} in the CMake.

The string utilities provide functionality in several general categories (not an exhaustive list):
\begin{itemize}
	\item \textbf{Case and character manipulation}: Conversion between upper- and lower-case strings, removal of characters, and character containment checks.
	
	\item \textbf{String parsing}: Splitting strings using delimiters and extracting portions of strings before, after, or between specified characters.
	
	\item \textbf{String validation}:  Checking whether strings contain digits, represent dice-roll notation, represent affirmative responses, or contain specified substrings and delimiters.
	
	\item \textbf{String transformation}: Shuffling, reversing, trimming, and otherwise transforming strings.
	
	\item \textbf{Type conversion}: Conversion between strings and vectors of characters or integer values.
	
	\item \textbf{Formatting}: Utilities for producing formatted strings, including centering text within a specified width.
	
	\item \textbf{Miscellaneous utilities}: Functions such as retrieving the current date and generating specialized string representations.
\end{itemize}

\subsubsection{Example}

The following demonstrates basic use of the string utilities:

\begin{lstlisting}[style=cppstyle]
#include "StringUtils.hpp"

std::string input = "  Hello World  ";

std::string upper = string_utils::toUpper(input);
std::string lower = string_utils::toLower(input);
std::string trimmed = string_utils::trim(input);
std::string reversed = string_utils::invertString(input);
\end{lstlisting}
The resulting strings are:
\begin{lstlisting}[style=cppstyle]
// input:    "  Hello World  "
// upper:    "  HELLO WORLD  "
// lower:    "  hello world  "
// trimmed:  "Hello World"
// reversed: "  dlroW olleH  "
\end{lstlisting}




\subsection{Vector Utilities}\label{section:vector_utils}

The MIA project provides a collection of general-purpose utilities for working with C++ vectors. The API for these utilities is defined in \texttt{VectorUtils.hpp}, where the available functions and their detailed behavior are documented. To use the vector utilities, include the \texttt{VectorUtils.hpp} header and link the application against the \texttt{Types\_UTIL} library in CMake.
















\subsection{Python Utilities}\label{section:python_utils}

The MIA project contains python utilities which allow loading and calling methods from python files from the c++ code. The python related utilities are contained in the \texttt{Python\_UTIL} library in CMake.

The \texttt{PythonModule.hpp} file provides a utility for importing and calling methods from a python file. To import a python module for use in c++ code, the python module (file) should typically be located in the same directory as the c++ code that is using it. The \texttt{PythonModule} is intended to be constructed with the name of the module (python file) and the file location of the calling file (this is used to determine where the python files are located).
\begin{lstlisting}[style=cppstyle]
// PythonModule(const std::string& moduleName, const std::string& callerFile);
PythonModule myModule("myPythonFile", __FILE__);
\end{lstlisting}
This class is intended to be called with \texttt{\_\_FILE\_\_} as the second parameter - which assumes the python files are located next to the associated c++ files. The \texttt{callerFile} parameter is used to determine the \textit{directory} of the python files, so this can be customized if the files are placed in a different directory (at the time of writing this, that is untested). See section \ref{section:python_file_location} for more information about the python file locations. 

After constructing the module, the python methods can simply be called using the appropriate \texttt{call()} overload:
\begin{lstlisting}[style=cppstyle]
myModule.call("pythonMethodName", <appropriate_parameters>);
\end{lstlisting}
The appropriate parameters that are fed into the \texttt{PythonModule::call(..)} method are parsed using a \texttt{PyObjectPtr toPython(..)} method. An overload of this method exists for each supported type, and a new overload may need added when adding python methods which do not use parameters of the existing supported types.


\subsubsection{Python Results}

The \texttt{PythonModule::call(..)} method returns a special \texttt{PythonResult} object. This object is used to store the return values of the python methods. This object is essentially a wrapper around all of the various supported types that python methods in the MIA project can return. If a non-supported type is returned, this object will need expanded to include new types. A failed call (a missing method, a python exception, or a conversion failure) does not throw; it returns an invalid result whose description is available through \texttt{getError()}, so callers should check \texttt{isValid()} before reading a value. The value is read through the \texttt{asInt()}, \texttt{asDouble()}, or \texttt{asString()} accessor matching the \texttt{getType()} of the result. The full result object API can be seen in the \texttt{PythonResult.hpp} file. Only module construction throws, when the module fails to import; a constructed module is always usable.



\subsubsection{Building Python Features}

The MIA project contains a CMake specific flag for building the python features. When building, if the python packages are located successfully, CMake will set the \texttt{BUILD\_PYTHON\_FEATURES} variable to true. This flag is intended to wrap features which require the python components so that they can be optionally built only when python libraries are available. The CMake for an app which requires the python files should therefore look like the following:
\begin{lstlisting}[style=makestyle]
# This app currently requires the python features to be available.
if(BUILD_PYTHON_FEATURES)
	# App-specific cmake here.
endif()
\end{lstlisting}


\subsubsection{Python File Location}\label{section:python_file_location}

When building the MIA project, the python install locations are set relative to the application install location (see \ref{section:platform_specific_path_configurations}). These paths are a separate python folder alongside the program files location:
\begin{lstlisting}[style=makestyle]
set(PYTHON_INSTALL_LOCATION "${APP_INSTALL_LOCATION}/python")
\end{lstlisting}
The \texttt{PythonModule} utility will attempt to locate the associated python modules for various cases. If the application is running from a system-installed location, it will look in the system python directory. Next, if a 'python' subdirectory of the resources folder next to the executable exists (\texttt{<exe\_dir>/resources/python}), this will be used (the release-build layout). Otherwise, the directory of the calling file (or custom path if \texttt{callerFile} is modified) will be used, since python files typically sit beside the code which uses them during development.







section{UI Input Utility}\label{section:ui_utils}

These features are all part of the \texttt{UI\_UTIL}.

The \texttt{UI\_UTIL} library contains the input event classes in \texttt{bin/utils/ui/}. The \texttt{Event} class wraps the plain string (intended to be expanded as needed) which a user interface reports for one user action (e.g. ``increment'' or ``stop''). The \texttt{EventQueue} class is a thread-safe queue which collects events from one or more producers for a single reader (handler); the reader drains it through \texttt{getAllEvents()} (non-blocking) or \texttt{waitForEvents()} (blocks until events arrive or the queue is closed). The \texttt{EventStorage} class is the batch of events returned by one read from the queue. Because all producers share one queue, the handler has a single place to wait for events to arrive, where it can then process them.

\subsection{Keybind Listener}

The \texttt{KeybindListener} class is a background task (see section \ref{section:background_task}) which watches keys and key combinations and pushes an event into an \texttt{EventQueue} when a binding completes. Bindings are read either from a \texttt{MIAConfig} (pairs whose keys start with a prefix, e.g. the pair \texttt{keybind\_a+b} = \texttt{increment}) or from a map of key sets to events. A binding fires when all of its keys are held at the same time, optionally together with a named modifier (\texttt{alt}, \texttt{ctrl}, or \texttt{shift}). A binding fires once per completion and does not repeat until one of its keys is released.

The listener also answers state queries: \texttt{isKeyHeld()} checks one key, and \texttt{isBindingHeld()} checks whether every key of a set (and its modifier, if any) is currently held.

\subsection{Single Key Listener}

The \texttt{SingleKeyListener} class is an older, simpler background task which watches one key and toggles an internal condition each time that key is pressed. The reading loop checks the condition through \texttt{isConditionMet()}; see section \ref{section:background_task} for an example of that pattern. This was first developed for the MIA Sequencer application (see chapter \ref{sequencer}) and later moved to the ui utility location as these types of features began expanding.

\subsection{Platform Details}

Both listeners detect presses no matter which window has focus. On Linux, detection polls the physical keyboard state through X11, and the \texttt{KeybindListener} only captures (swallows) its keys when \texttt{captureGlobally} is set; on Windows, the key states are polled with \texttt{GetAsyncKeyState} and never captured. The key listeners are optional: they build on Windows, and on Linux only when the libxdo and X11 development libraries are found (see section \ref{section:optional_features}); the \texttt{Event} classes always build.

\subsubsection{Defining Bindings}

There are two ways to give the listener its bindings.

\paragraph{From a config.} The first constructor reads a KEY\_VALUE config (see section \ref{section:configuration_system}) and binds every pair whose key starts with the keybind prefix (\texttt{keybind\_} by default). The tokens after the prefix name the keys; the pair's value names the event which is pushed when the binding completes:
\begin{lstlisting}[style=makestyle]
# The character key w pushes the "up" event.
keybind_w=up

# '+' joins keys into a combination, held at the same time.
keybind_a+b=both

# The tokens 'alt', 'ctrl', and 'shift' are modifiers; any number of
# them may be combined with any number of character keys.
keybind_alt+t=alt_t
keybind_alt+shift+s=alt_shift_s
\end{lstlisting}
Pairs which do not start with the prefix, or whose tokens are not characters or modifier names, are skipped (reported to standard error when verbose mode is set). The config must be initialized before it is passed in; the constructor only reads from it.

\paragraph{From a map.} The second constructor takes a \texttt{KeybindEventMap}, which maps a \texttt{Keybind} (a set of characters plus a \texttt{Modifier}) directly to the event it pushes. This is for bindings built in code rather than read from a file.

\subsubsection{Using the Listener}

A listener follows the background task lifecycle (see section \ref{section:background_task}): construct it, call \texttt{initialize()}, then \texttt{start()}; call \texttt{stop()} when finished. The completed-binding events are read from the shared queue, either by polling \texttt{getAllEvents()} or by blocking in \texttt{waitForEvents()}:
\begin{lstlisting}[style=cppstyle]
ui::EventQueue queue;
ui::KeybindListener listener(config, queue);
listener.initialize();
listener.start();

while (listener.isRunning())
{
    ui::EventStorage storage = queue.waitForEvents();
    for (const ui::Event& event : storage.getEvents())
        handleEvent(event.getName());
}

listener.stop();
\end{lstlisting}

For movement-style input, which lasts as long as the key is held, the reading loop checks the held state each frame instead of waiting for a completion:
\begin{lstlisting}[style=cppstyle]
while (listener.isKeyHeld('w'))
{
    moveForward();
    timing::sleepMilliseconds(10);
}

// A chord: true only while both keys are held together.
while (listener.isBindingHeld({'a', 'd'}))
{
    strafe();
    timing::sleepMilliseconds(10);
}
\end{lstlisting}
Only keys which appear in a binding are tracked; a key never bound by the listener always reads as not held. Single presses (including ones with a required modifier, e.g. \texttt{ctrl+s}) are simpler as completion bindings, since the queue push already provides the edge.

The constructor parameters (the config key prefix, the poll interval, and global key capture) are documented in \texttt{KeybindListener.hpp}; the class comment there also covers the capture behavior and the re-arm semantics in full.
















\section{Audio Utility}
\label{section:audio}

These features are all part of the \texttt{Audio\_UTIL}. The \texttt{Audio\_UTIL} utility plays system sounds and audio files. Its API lives in \texttt{bin/utils/audio/} across four headers: \texttt{SystemSounds.hpp} for system-specific sounds (e.g., beeps), \texttt{SoundsFromFile.hpp} for one-shot file playback (primarily used as proof of concept while developing the \texttt{AudioPlayer}), \texttt{AudioPlayer.hpp} for managed playback, and \texttt{Playlist.hpp} for the ordered track list with shuffle memory which backs the player. This section covers these various audio interfaces.

\subsection{System Sounds}

The \texttt{system\_sounds} namespace emits simple system-level cues. 

\subsubsection{System Beeps}

The \texttt{beep()} method plays one beep at a given frequency for a given duration in milliseconds:
\begin{lstlisting}[style=cppstyle]
// The Zelda area-unlock jingle.
system_sounds::beep(784, 100);
system_sounds::beep(740, 100);
system_sounds::beep(1047, 500);
\end{lstlisting}
On Windows this wraps the \texttt{Beep} API. On Linux the frequency and duration are not yet supported; however, the method opens \texttt{/dev/console} and writes a terminal bell character to create a beep.

\subsection{Playing a Sound File Once}

The \texttt{SoundsFromFile.hpp} free functions play a file and forget it. \texttt{playSoundFromFile()} blocks the calling thread until the file finishes; \texttt{playSoundFromFileAsync()} returns immediately and plays on a background thread, with \texttt{stopSound()} stopping whatever is playing:
\begin{lstlisting}[style=cppstyle]
#include "SoundsFromFile.hpp"

// Blocks until the sound finishes.
audio::playSoundFromFile("alarm.mp3");

// Plays in the background; stop after 5 seconds.
audio::playSoundFromFileAsync("music.mp3");
audio::stopSound(0);
\end{lstlisting}
\texttt{stopSound()} takes an optional fade-out length in milliseconds which ramps the volume down before stopping. The fade uses MCI volume commands which only support mp3 files, so a fade requested for a wav file stops immediately and prints a warning. These functions track a single playing sound globally; for anything more, use the \texttt{AudioPlayer} class.

\begin{note}
The entire free-function API was essentially a proof of concept exploration of playing audio files through code. After getting it to work, it was decided that a class object is a much cleaner way of approaching the interface and the \texttt{AudioPlayer} (see section \ref{section:audio_player}) was then developed using much of the same code with some added features.
\end{note}

\subsection{The AudioPlayer Class}\label{section:audio_player}

The \texttt{AudioPlayer} class is a background task (see section \ref{section:background_task}) which plays audio through its own worker thread, so the calling thread stays free. It reports state through \texttt{isAudioPlaying()} and is controlled with \texttt{startAudio()}, \texttt{restartAudio()}, and \texttt{stopAudio()}:
\begin{lstlisting}[style=cppstyle]
#include "AudioPlayer.hpp"

audio::AudioPlayer player("alarm.mp3");
player.startAudio();

// ...the audio plays in the background...

// Stop immediately, or fade out over three seconds.
player.stopAudio(0);
player.stopAudio(3000);
\end{lstlisting}
\texttt{startAudio()} does nothing when the audio is already playing, and \texttt{restartAudio()} starts the current track over. A fade-out length only takes effect on mp3 files on Windows for the same MCI reason as above. On Linux, the fade works for all supported file types. The fade-out runs asynchronously in the player's worker thread, so \texttt{stopAudio()} returns before the fade completes and the player reports \texttt{isAudioPlaying()} as false once the fade has finished.

The playback volume can be set with \texttt{setVolume()}, which takes a value from 0 to 100 (larger values are clamped), and read back with \texttt{getVolume()}:
\begin{lstlisting}[style=cppstyle]
audio::AudioPlayer player("music.mp3");
player.setVolume(50);
player.startAudio();

// ...after a few seconds...
player.setVolume(100); // Applies immediately, mid-play.
\end{lstlisting}
The volume applies immediately to the current playback and persists across tracks, restarts, and playlists until it is changed again.

The player always holds a playlist. Constructing it from a file name or calling \texttt{setAudioFile()} is just a one-track playlist; the playlist methods manage the full list instead:
\begin{lstlisting}[style=cppstyle]
audio::AudioPlayer player;
player.setPlaylist({"track1.mp3", "track2.mp3", "track3.wav"});
player.addToPlaylist("bonus.mp3");
player.startAudio();
\end{lstlisting}
When a track finishes, the player moves to the next one and keeps going until the playlist is exhausted; only then does it stop on its own. Three independent settings shape that behavior:
\begin{itemize}\itemsep0em
	\item \texttt{setRepeatAudioTrack(bool)} repeats the current track instead of advancing.
	\item \texttt{setPlaylistLoop(bool)} wraps the playlist back to the first track after the last one ends.
	\item \texttt{setPlaylistShuffle(bool)} picks the next track at random rather than in order. The pick is never the track being left, so shuffle never plays the same track twice in a row.
\end{itemize}
The playlist state survives edits made during playback: \texttt{setPlaylist()}, \texttt{addToPlaylist()}, and \texttt{setAudioFile()} take effect on the next track rather than interrupting the current one. \texttt{nextAudio()} and \texttt{previousAudio()} skip manually mid-play. The player keeps a memory of previously played tracks, so \texttt{previousAudio()} returns to the track which played before the current one, shuffle or not; with no memory yet, it prints a warning and leaves playback running. When the audio file is changed mid-play, the currently playing track is no longer part of the playlist, so \texttt{nextAudio()} plays the track at the playlist cursor (the newly set file) instead of advancing past it.

\subsection{Supported File Types}

Currently, the following types are supported, with a few limitations:
\begin{itemize}
	\item \texttt{MP3}: The MP3 file type has full support for all features on all supported platforms.
	\item \texttt{WAV}: The WAV file type is supported on Windows and supports all features except the optional fade-out when stopping the audio. The limitation exists because the current Windows MCI implementation uses the \texttt{setaudio ... volume to} command to gradually reduce volume during a fade-out. MCI supports volume control for MP3 playback, but the volume command is not supported for WAV playback. On Linux, the fade uses the player's own volume control, so there is no such limitation.
\end{itemize}

\subsection{Platform Details}

\begin{itemize}\itemsep0em
	\item \textbf{Windows:} All playback goes through the Windows Media Control Interface (MCI, from \texttt{winmm}). Supported file types are mp3 and wav; the supported set is checked through \texttt{audio::isASupportedType()} before any playback starts.
	\item \textbf{Linux:} All file playback goes through libVLC (from \texttt{libvlc}). Only mp3 is in the supported set so far. The player creates a fresh VLC media player for each playback and drives the volume through libVLC's audio API, which supports fades for every supported file type. The \texttt{libvlc-dev} package is required to build; the setup script installs it. The system beep works through \texttt{/dev/console}.
\end{itemize}
